\documentclass[10pt,letterpaper]{article}
\usepackage[margin=1in,headheight=14pt]{geometry}
\usepackage[T1]{fontenc}
\usepackage{lmodern}
\usepackage{amsmath,amssymb,amsthm,mathtools}
\usepackage{microtype,booktabs,xcolor,fancyhdr}
\definecolor{ink}{HTML}{19354A}
\usepackage[colorlinks=true,linkcolor=ink,urlcolor=ink,citecolor=ink]{hyperref}
\hypersetup{pdftitle={Rainbow progressions in small intervals and groups},pdfauthor={Akhil Ravuri}}
\newtheorem{proposition}{Proposition}
\newcommand{\aw}{\operatorname{aw}}
\pagestyle{fancy}\fancyhf{}
\fancyhead[L]{\small\textsc{Rainbow progressions}}
\fancyhead[R]{\small Akhil Ravuri}\fancyfoot[C]{\thepage}
\renewcommand{\headrulewidth}{0.3pt}
\setlength{\parskip}{0.2em}\setlength{\emergencystretch}{1em}
\title{\vspace{-1.5em}\LARGE\bfseries Rainbow progressions\\in small intervals and groups}
\author{Akhil Ravuri}\date{}
\begin{document}
\maketitle\vspace{-1.5em}
\begin{abstract}
We record exhaustive anti-van der Waerden threshold calculations for intervals
and cyclic groups of orders three through nine, and three small product
groups. Definitions follow the cited literature. The searches reproduce
selected known benchmarks and illustrate why the ambient group matters.
These are bounded finite computations, without a novelty claim.
\end{abstract}

\section{Definitions and the threshold convention}
A nondegenerate three-term arithmetic progression is a set of three distinct
elements $\{a,a+d,a+2d\}$. For an interval we require $d>0$; in a finite
abelian group we use its addition and remove duplicate sets. Reversal gives
the same set. A progression is rainbow if its three elements have distinct
colors. The interval $[n]=\{1,\ldots,n\}$ is represented in code by
$\{0,\ldots,n-1\}$, which is equivalent by translation.

The threshold $\aw(S,3)$ is the smallest exact number of colors forcing a
rainbow progression, with an exact coloring required to use every color.
Definitions for intervals and cyclic groups follow Butler et al.~\cite{BE};
the finite-abelian-group formulation follows Young~\cite{Y}.

The search first maximizes the number of classes in an avoiding partition and
then adds one. This translation is valid because merging classes cannot
create a rainbow progression: a maximum avoiding coloring can be merged to
any smaller positive color count. If no nondegenerate progression exists, we
use the convention $\aw(S,3)=|S|+1$. This lies beyond feasible palette size;
it does not assert a coloring with more than $|S|$ colors.

\section{Exact interval and cyclic results}
Each domain is searched separately over every partition, ignoring color names
but keeping its elements fixed. Cyclic addition may wrap around.
\begin{center}
\begin{tabular}{rrrr}
\toprule
$n$ & Partitions per domain, $B_n$ & $\aw([n],3)$ & $\aw(\mathbb Z_n,3)$\\
\midrule
3 & 5 & 3 & 3\\
4 & 15 & 4 & 3\\
5 & 52 & 4 & 3\\
6 & 203 & 4 & 4\\
7 & 877 & 4 & 3\\
8 & 4140 & 5 & 3\\
9 & 21147 & 4 & 4\\
\bottomrule
\end{tabular}
\end{center}
For $3\le n\le8$, the interval values agree with~\cite[Theorem 1.1]{BE}.
The threshold is not assumed monotone: the interval value drops from five
at $n=8$ to four at $n=9$. These entries concern exact-color forcing, and
should not be replaced by a monotonicity assumption about interval inclusion.

\section{Products of cyclic groups}
For $G=\mathbb Z_{m_1}\times\cdots\times\mathbb Z_{m_t}$, all coordinates are
added modulo their own moduli. The generator enumerates every pair $(a,d)$
and keeps the resulting three-element sets exactly once.
\begin{center}
\begin{tabular}{lrrr}
\toprule
Group & Partitions & Distinct 3-AP sets & $\aw(G,3)$\\
\midrule
$\mathbb Z_2\times\mathbb Z_2$ & 15 & 0 & 5\\
$\mathbb Z_2\times\mathbb Z_4$ & 4140 & 16 & 5\\
$\mathbb Z_2^3$ & 4140 & 0 & 9\\
\bottomrule
\end{tabular}
\end{center}
\begin{proposition}
For every integer $t\ge1$, $\mathbb Z_2^t$ has no nondegenerate three-term
arithmetic progression. Under the stated convention,
$\aw(\mathbb Z_2^t,3)=2^t+1$.
\end{proposition}
\begin{proof}
Every $d\in\mathbb Z_2^t$ satisfies $2d=0$, so $a+2d=a$. Thus the first and
third terms of any proposed progression agree. No three-element progression
exists, and the threshold follows from the no-progression convention.
\end{proof}
This elementary known observation explains the zero-progression entries.
Groups of order eight therefore need not behave alike: $\mathbb Z_8$,
$\mathbb Z_2\times\mathbb Z_4$, and $\mathbb Z_2^3$ give thresholds three,
five, and nine respectively. The product examples are direct enumerations,
not implementations of Young's general group formulas or unitary colorings.

\section{Coverage, checks, and reproduction}
The implementation enumerates all restricted-growth strings of each domain
size. Bell-number coverage is checked using an independent Stirling-number
recurrence. Interval progressions are checked against all triples satisfying
$a+c=2b$; cyclic progressions are checked against triples with some ordering
satisfying that equation modulo $n$. Product-group tests independently check
the coordinate equations. These checks audit the finite implementation;
the general exponent-two observation above has its own proof.

From the repository root, with Python 3.11+ and its standard library, run
\begin{quote}\small\ttfamily
python -m experiments.rainbow\_progressions\par
python -m experiments.reproduce --check results/data/small\_cases.json
\end{quote}
The reference JSON saves counts and avoiding witnesses. Progression generation
is in \path{src/arithmetic_progressions.py}. Partition enumeration grows
rapidly; the recorded range is deliberately small and gives no fast general
threshold algorithm. The
\href{https://github.com/Akhil-rav/discrete-math-lab/blob/main/notes/literature/source_records.md}{source records}
document the consulted versions, locators, and PDF internal-date discrepancies.
No claim of a new theorem or currently open problem is made.

\begin{thebibliography}{2}
\bibitem{BE} S. Butler et al., \emph{Rainbow arithmetic progressions},
\href{https://arxiv.org/abs/1404.7232v2}{arXiv:1404.7232v2}.
Definitions, Theorem 1.1, Remark 1.4, and cyclic discussion checked.
\bibitem{Y} M. Young, \emph{Rainbow Arithmetic Progressions in Finite Abelian Groups},
\href{https://arxiv.org/abs/1603.08153v1}{arXiv:1603.08153v1}.
Group definitions and Theorem 3 statement checked; general formulas are not
implemented here.
\end{thebibliography}
\end{document}
